% © 2026 Ing. David Jaroš — CC BY-NC-ND 4.0
%
% This work is licensed under a Creative Commons Attribution-NonCommercial-NoDerivatives
% 4.0 International License (CC BY-NC-ND 4.0).
%
% License History: Earlier drafts (up to v0.3) were released under CC BY 4.0.
% From v0.4 onward, all material is released under CC BY-NC-ND 4.0 to protect
% the integrity of the theoretical work during ongoing academic development.
%
% See LICENSE.md for full license text.

\ifdefined\INCLUDEMODE
\else
  \documentclass[12pt,a4paper]{article}
  \usepackage[a4paper, margin=2.5cm]{geometry}
  \usepackage{amsmath, amssymb, amsthm}
  \usepackage{mathtools}
  \usepackage{hyperref}
  \usepackage{xcolor}
  \usepackage{mdframed}
  \usepackage{booktabs}
  \usepackage{longtable}
  \usepackage{tcolorbox}

  \newtheorem{theorem}{Theorem}[section]
  \newtheorem{lemma}[theorem]{Lemma}
  \newtheorem{proposition}[theorem]{Proposition}
  \newtheorem{corollary}[theorem]{Corollary}
  \theoremstyle{definition}
  \newtheorem{definition}[theorem]{Definition}
  \theoremstyle{remark}
  \newtheorem{remark}[theorem]{Remark}

  \newmdenv[backgroundcolor=yellow!20, linecolor=orange, linewidth=1.5pt]{gapbox}
  \newmdenv[backgroundcolor=green!15, linecolor=green!60!black, linewidth=1.5pt]{resultbox}
  \newmdenv[backgroundcolor=red!8, linecolor=red!60, linewidth=1.5pt]{warnbox}
  \newmdenv[backgroundcolor=blue!8, linecolor=blue!50, linewidth=1.5pt]{insightbox}

  \title{\textbf{Best Route to $\alpha$ from UBT: Prime Attractor on $S^1_\psi$}\\
         \large Alpha Route Status Review — Formal Derivation}
  \author{Ing.\ David Jaro\v{s}}
  \date{2026-04-28}

  \begin{document}
  \maketitle
\fi

%──────────────────────────────────────────────────────────────────────────────
\section{Purpose and Status}
\label{sec:ar:purpose}
%──────────────────────────────────────────────────────────────────────────────

This document presents the \textbf{conditional structural route to
$\alpha^{-1}_{\mathrm{bare}} = 137$}, blocked by Gap G137-B.
No first-principles derivation of the physical value
$\alpha^{-1} = 137.036$ is claimed here.

\begin{warnbox}
\textbf{Canonical status sync (2026-05-10)}:
alpha is \textbf{NOT derived};
$\alpha^{-1}_{\mathrm{bare}} = 137$ is \textbf{CONDITIONAL ONLY};
Gap \textbf{G137-B} remains open.
Legacy references in this file to older G137-B framing are historical and
do not override \texttt{canonical/alpha/ALPHA\_MASTER\_STATUS.md}.
\end{warnbox}

\begin{tcolorbox}[colback=blue!4!white, colframe=blue!50!black,
                  title=Mission Classification]
\begin{tabular}{ll}
Route name & Prime Attractor on $S^1_\psi$ \\
Target & $\alpha^{-1}_{\mathrm{bare}} = 137$ (bare, UV value) \\
Proved steps & 6/7 (Steps 1--4, 6, and prime-stability part of Step 7) \\
Conditional step & Derive $B$ from $S[\Theta]$ without circular inputs (Gap G137-B) \\
Circular inputs & Zero (in bare-value claim only) \\
Overall status & \textbf{CONDITIONAL on Gap G137-B} \\
Hard rule & No constant chosen to match $\alpha$ \\
\end{tabular}
\end{tcolorbox}

\begin{warnbox}
\textbf{Hard rule}: No numerical parameter in this document may be tuned to
match~$\alpha$.  Every constant is either derived from UBT axioms or explicitly
marked as a Motivated Conjecture~[MC].  All circular inputs (which use $\alpha$
or $m_e$) are marked~[CIRC] and excluded from the first-principles claim.
\end{warnbox}

%──────────────────────────────────────────────────────────────────────────────
\section{Derivation Chain Overview}
\label{sec:ar:chain}
%──────────────────────────────────────────────────────────────────────────────

The derivation proceeds in seven steps:
\begin{equation}
  \underbrace{\mathcal{B}=\mathbb{C}\otimes\mathbb{H}}_{\text{[L0]}}
  \;\to\;
  \underbrace{N_{\mathrm{eff}}=12}_{\text{[L0]}}
  \;\to\;
  \underbrace{B_0=8\pi}_{\text{[L1]}}
  \;\xrightarrow{\text{Gap G137-B}}\;
  \underbrace{B\approx 46.298\ \text{from}\ S[\Theta]}_{\text{[OPEN]}}
  \;\to\;
  \underbrace{V_{\mathrm{eff}}(n)}_{\text{[L1]}}
  \;\to\;
  \underbrace{n^*=137}_{\text{[L1]}}
  \;\to\;
  \underbrace{\alpha^{-1}_{\mathrm{bare}}=137.}_{\text{[COND]}}
  \label{eq:ar:chain}
\end{equation}

Steps marked [L0] are axiomatic consequences.
Steps marked [L1] are one-loop field-theory derivations.
The single conditional step is marked [MC] and depends on proving $k = 1$.

%──────────────────────────────────────────────────────────────────────────────
\section{Step 1: Biquaternion Algebra Axiom}
\label{sec:ar:algebra}
%──────────────────────────────────────────────────────────────────────────────

\begin{definition}[UBT fundamental algebra]
\label{def:ar:algebra}
The canonical field algebra of UBT is:
\begin{equation}
  \mathcal{B} = \mathbb{C}\otimes_{\mathbb{R}}\mathbb{H}
  \;\cong\; M_2(\mathbb{C}),
  \label{eq:ar:B}
\end{equation}
with $\dim_{\mathbb{R}}\mathcal{B} = 8$ and imaginary quaternion subspace
$\mathrm{Im}\,\mathbb{H} \cong \mathbb{R}^3$, $\dim_{\mathbb{R}}(\mathrm{Im}\,\mathbb{H}) = 3$.
\end{definition}

\begin{resultbox}
Status: CLEAN [L0] — canonical axiom, no reference to $\alpha$ or $m_e$.\\
Source: \texttt{canonical/fields/biquaternion\_algebra.tex}
\end{resultbox}

%──────────────────────────────────────────────────────────────────────────────
\section{Step 2: Complex Time and $S^1_\psi$ Compactification}
\label{sec:ar:compactification}
%──────────────────────────────────────────────────────────────────────────────

\begin{definition}[$S^1_\psi$ and Dirac quantisation]
\label{def:ar:S1psi}
The imaginary component $\psi$ of complex time $\tau = t + i\psi$ is
compactified on a circle of radius $R_\psi$:
\begin{equation}
  \psi \;\sim\; \psi + 2\pi R_\psi,
  \label{eq:ar:Rpsi_circle}
\end{equation}
yielding a Kaluza-Klein tower of winding modes $n \in \mathbb{Z}$.
The Dirac quantisation condition on the charged $\Theta$-field:
\begin{equation}
  e^{iq\oint_\psi A_\psi\,d\psi} = 1
  \label{eq:ar:dirac}
\end{equation}
follows from unitarity and gauge consistency.
\end{definition}

\begin{resultbox}
Status: CLEAN [L0] — follows from $S[\Theta]$ gauge symmetry.\\
Source: \texttt{canonical/appendices/appendix\_alpha\_geometry.tex \S1}
\end{resultbox}

%──────────────────────────────────────────────────────────────────────────────
\section{Step 3: $N_{\mathrm{eff}} = 12$ Charged Modes}
\label{sec:ar:Neff}
%──────────────────────────────────────────────────────────────────────────────

\begin{theorem}[$N_{\mathrm{eff}} = 12$]
\label{thm:ar:Neff}
The effective number of distinct charged modes of the biquaternion field on $S^1_\psi$
contributing to the one-loop vacuum polarisation is:
\begin{equation}
  N_{\mathrm{eff}}
  = \underbrace{3}_{N_{\mathrm{phases}}} \times
    \underbrace{2}_{N_{\mathrm{helicity}}} \times
    \underbrace{2}_{N_{\mathrm{charge}}}
  = 12,
  \label{eq:ar:Neff}
\end{equation}
where:
\begin{itemize}
  \item $N_{\mathrm{phases}} = \dim_{\mathbb{R}}(\mathrm{Im}\,\mathbb{H}) = 3$
        (three independent quaternion imaginary directions $\mathbf{i,j,k}$,
        each carrying an independent $U(1)$ phase);
  \item $N_{\mathrm{helicity}} = 2$ (left/right helicity of complexified $\Theta$);
  \item $N_{\mathrm{charge}} = 2$ (particle + antiparticle: complex structure of $\mathcal{B}$).
\end{itemize}
\end{theorem}

\begin{proof}
The biquaternion field $\Theta \in \mathcal{B} = \mathbb{C}\otimes\mathbb{H}$.
The three imaginary quaternion directions are independent phase directions;
each supports left and right helicity components of the complexified field;
and the overall complex structure provides charge-conjugation doubling.
All three factors are independent and multiply by construction.
\end{proof}

\begin{resultbox}
Status: CLEAN [L0] — zero-free-parameter algebraic theorem.\\
Circularity: None.  Verified: \texttt{canonical/n\_eff/step3\_N\_eff\_result.tex}
\end{resultbox}

%──────────────────────────────────────────────────────────────────────────────
\section{Step 4: One-Loop Coefficient $B_0 = 8\pi$}
\label{sec:ar:B0}
%──────────────────────────────────────────────────────────────────────────────

\begin{theorem}[$B_0 = 8\pi$ from one-loop vacuum polarisation]
\label{thm:ar:B0}
The one-loop vacuum polarisation of $N_{\mathrm{eff}} = 12$ charged scalar modes
on the compact imaginary-time circle $S^1_\psi$ gives the baseline coefficient:
\begin{equation}
  B_0 = \frac{2\pi N_{\mathrm{eff}}}{3}
      = \frac{24\pi}{3}
      = 8\pi
      \approx 25.133.
  \label{eq:ar:B0}
\end{equation}
\end{theorem}

\begin{proof}
The one-loop correction to the effective potential for a $U(1)$ winding mode
from $N_{\mathrm{eff}}$ charged scalar modes in one (compact) dimension is
$-B_0 \ln n$ with $B_0 = 2\pi N_{\mathrm{eff}}/3$.  The factor $1/3$ arises
from the standard one-loop vacuum polarisation integral in one dimension;
the $2\pi$ integrates the torus measure.  In the QED limit $N_{\mathrm{eff}}=1$,
one recovers $B_0 = 2\pi/3$ (standard one-loop QED result, consistent check).
\end{proof}

\begin{resultbox}
Status: CLEAN [L1] — standard one-loop field theory; no free parameters.\\
Source: \texttt{canonical/n\_eff/step2\_vacuum\_polarization.tex}
\end{resultbox}

%──────────────────────────────────────────────────────────────────────────────
\section{Step 5 (Conditional): $B_{\mathrm{base}} = N_{\mathrm{eff}}^{3/2}$}
\label{sec:ar:Bbase}
%──────────────────────────────────────────────────────────────────────────────

This is the \textbf{only conditional step} in the chain.

\begin{gapbox}
\textbf{Gap G137-B (the unique blocking gap)}: Prove that the Kac-Moody level of
the WZW-type current algebra describing the biquaternion field on the
$\psi$-torus is $k = 1$.
\end{gapbox}

\begin{proposition}[$B_{\mathrm{base}} = N_{\mathrm{eff}}^{3/2}$, given $k = 1$]
\label{prop:ar:Bbase}
Assume the biquaternion field theory on $T^2_\tau$ is a 2D CFT with an
$\mathfrak{su}(2)$ Kac-Moody algebra at level $k$.  The full one-loop coefficient
for the $S^1_\psi$ winding-mode potential is:
\begin{equation}
  B_{\mathrm{base}} = N_{\mathrm{eff}}^{3/2} \cdot k^{1/2}.
  \label{eq:ar:Bbase_formula}
\end{equation}
If $k = 1$:
\begin{equation}
  B_{\mathrm{base}} = N_{\mathrm{eff}}^{3/2} = 12^{3/2} = 12\sqrt{12} \approx 41.57.
  \label{eq:ar:Bbase_k1}
\end{equation}
\end{proposition}

\begin{insightbox}
\textbf{Structural motivation for $k = 1$} (not a proof):

\textbf{(i) Modular partition function}:
The partition function of the UBT biquaternion field on the torus is
$\hat{Z}(\tau) = \vartheta_3^3(\tau)$, which transforms with modular weight
$3/2$ under $SL(2,\mathbb{Z})$ (Gap G8, closed).  The weight $3/2$ equals
the exponent in $N_{\mathrm{eff}}^{3/2}$, suggesting a common algebraic origin
in the three independent imaginary quaternion directions.

\textbf{(ii) New analysis (Alpha route review 2026-04-28)}:
$\vartheta_3^3(\tau)$ is the partition function of three compact free bosons
at the self-dual radius, which is \emph{also} the three-fold product of
$SU(2)_1$ WZW partition functions (these coincide at the self-dual point).
The $SU(2)_1$ WZW model has central charge $c = 3/(1+2) = 1$ per factor;
three factors give $c = 3$.  This is consistent with $k = 1$ per imaginary
quaternion direction.  However, this observation does not yet constitute a proof:
$\vartheta_3^3$ is also consistent with three free (non-WZW) bosons at
$c = 3$, $k \to \infty$.  A modular bootstrap computation is required to
distinguish between these possibilities.

\textbf{(iii) Absence of Chern-Simons term}:
The action $S[\Theta]$ is parity-symmetric; no Chern-Simons term is present.
In 2D CFT, this is consistent with the minimal level $k = 1$ (free-field fixed
point with no topological term).

\medskip
\textbf{Status of Gap G137-B}: All 27+ approaches exhausted except the modular
bootstrap.  The modular bootstrap (crossing symmetry on the 4-point function
on $T^2$) is the one remaining untested route.  Time-box: 4 weeks from 2026-04-28.
\end{insightbox}

\begin{resultbox}
Status: [MC] Motivated Conjecture — conditional on $k = 1$ proof.\\
If proved: CLEAN [L1].  No circular input.\\
Source: \texttt{canonical/interactions/B\_base\_derivation\_complete.tex}
\end{resultbox}

%──────────────────────────────────────────────────────────────────────────────
\section{Step 6: Effective Potential $V_{\mathrm{eff}}(n)$}
\label{sec:ar:Veff}
%──────────────────────────────────────────────────────────────────────────────

\begin{theorem}[One-loop effective potential]
\label{thm:ar:Veff}
The one-loop effective potential for winding mode $n \geq 1$ on $S^1_\psi$ is:
\begin{equation}
  V_{\mathrm{eff}}(n) = n^2 - B_{\mathrm{base}}\, n\ln n + \mathrm{const},
  \qquad B = B_{\mathrm{base}}.
  \label{eq:ar:Veff}
\end{equation}
The $n^2$ term is the classical KK winding energy; $-B\, n\ln n$ is the
one-loop correction from $N_{\mathrm{eff}}$ charged modes.
\end{theorem}

\begin{resultbox}
Status: CLEAN [L1] given $B_{\mathrm{base}}$.  Functional form is standard
one-loop; only coefficient tracks Gap G137-B.\\
Source: \texttt{canonical/appendices/appendix\_alpha\_geometry.tex \S3}
\end{resultbox}

%──────────────────────────────────────────────────────────────────────────────
\section{Step 7: Prime Attractor $n^* = 137$}
\label{sec:ar:nstar}
%──────────────────────────────────────────────────────────────────────────────

\begin{theorem}[Stationary winding number]
\label{thm:ar:nstar_stationary}
The continuous minimum of $V_{\mathrm{eff}}(n) = n^2 - B\,n\ln n$ satisfies:
\begin{equation}
  \frac{\partial V_{\mathrm{eff}}}{\partial n}\bigg|_{n^*} = 0
  \;\;\Longrightarrow\;\;
  2n^* = B_{\mathrm{base}}\bigl(\ln n^* + 1\bigr).
  \label{eq:ar:nstar_formula}
\end{equation}
Solving numerically for $B_{\mathrm{base}} \approx 41.57$, one obtains
$n^*_{\mathrm{continuous}} \approx 120.4$.
For the corrected coefficient $B_{\mathrm{full}} = B_{\mathrm{base}} \times R \approx 46.31$,
the equation $2n = 46.31(\ln n + 1)$ solves to $n^*_{\mathrm{continuous}} \approx 137.1$.
The prime minimiser of $V_{\mathrm{eff}}(p)$ over primes $p$ is:
\begin{equation}
  n^*_{\mathrm{prime}} = \operatorname{arg\,min}_{p\,\mathrm{prime}} V_{\mathrm{eff}}(p)
  = 137 \quad \text{for } B_{\mathrm{full}} \approx 46.31.
  \label{eq:ar:nstar_prime}
\end{equation}
\end{theorem}

\begin{remark}[Role of the correction factor $R$]
\label{rem:ar:Rfactor}
The prime minimiser of $V_{\mathrm{eff}}(p)$ for $B = B_{\mathrm{base}} \approx 41.57$
gives $n^*_{\mathrm{prime}} = 127$ (the prime nearest the continuous minimum).
With the correction $B_{\mathrm{full}} = B_{\mathrm{base}} \times R \approx 46.31$
($R \approx 1.114$), the prime minimiser shifts to $n^*_{\mathrm{prime}} = 137$.

The factor $R \approx 1.114$ accounts for higher-loop or finite-volume corrections
to the one-loop effective potential.  It is a \emph{motivated conjecture} (no
rigorous derivation); see the stress test in \S\ref{sec:ar:stresstests}.
\end{remark}

\begin{theorem}[Prime stability]
\label{thm:ar:prime_stability}
A winding vacuum at $n^*$ is stable under $\pi_1(S^1_\psi)$-deformations if and
only if $n^*$ is prime.
\end{theorem}

\begin{proof}
If $n^* = a \cdot b$ (composite, $a,b > 1$), sub-harmonic modes at winding numbers
$a$ and $b$ are topologically allowed (both lie in $\pi_1(S^1_\psi) = \mathbb{Z}$)
and are energetically below $n^*$ in pairs.  A small perturbation can split the
$n^*$-winding into $a$- and $b$-windings.  A prime winding number has no such
sub-harmonic decomposition and is therefore topologically stable under all
$\pi_1$-deformations.  The integer 137 is prime by direct verification.
\end{proof}

\begin{resultbox}
Status of prime stability: CLEAN [L1] — proved independently of $B_{\mathrm{base}}$.\\
Status of $n^* = 137$: CONDITIONAL [L1] — depends on $B_{\mathrm{base}}$ (Step 5).\\
Source: \texttt{canonical/appendices/appendix\_alpha\_geometry.tex \S4}
\end{resultbox}

%──────────────────────────────────────────────────────────────────────────────
\section{Step 8: Identification $\alpha^{-1}_{\mathrm{bare}} = n^*$}
\label{sec:ar:identification}
%──────────────────────────────────────────────────────────────────────────────

\begin{proposition}[$\alpha^{-1}_{\mathrm{bare}} = 137$]
\label{prop:ar:alpha}
The bare fine-structure constant (UV, unrenormalised value before QED running)
is identified with the prime vacuum winding number via the Dirac quantisation
condition~\eqref{eq:ar:dirac}:
\begin{equation}
  \alpha^{-1}_{\mathrm{bare}} = n^* = 137.
  \label{eq:ar:alpha_bare}
\end{equation}
The identification arises because the winding number $n^*$ sets the minimal
charge quantum on $S^1_\psi$, which (in natural units at the T-duality self-dual
point $R_\psi = R_t$) equals the electron charge $e$, giving
$\alpha = e^2/(4\pi) = 1/n^*$.
\end{proposition}

\begin{warnbox}
\textbf{What this derivation does NOT claim}:
\begin{enumerate}
  \item The physical value $\alpha^{-1} = 137.036$ is \textbf{not derived}:
        the extra $\delta \approx 0.036$
        is the two-loop QED renormalisation correction, which within UBT requires
        the UV cutoff scale~$\Lambda$ and the electron mass~$m_e$ as inputs —
        both currently circular (Gap A9).
  \item The physical value of $R_\psi$ in SI units: requires $m_e$ as input
        (Gap A10 — semi-empirical).
  \item The value $k = 1$: Gap G137-B remains unproved.
\end{enumerate}
\end{warnbox}

%──────────────────────────────────────────────────────────────────────────────
\section{Stress Tests for Hidden Fitting}
\label{sec:ar:stresstests}
%──────────────────────────────────────────────────────────────────────────────

Per the mission hard rules, every numerical result is stress-tested.

\subsection{Non-circularity of $N_{\mathrm{eff}} = 12$}

The derivation would be circular if $N_{\mathrm{eff}} = 12$ were chosen to
reproduce $\alpha^{-1} = 137$.  Testing other gauge-group sizes
(computed using the correct $V_{\mathrm{eff}}(p) = p^2 - B_{\mathrm{full}}\,p\ln p$):

\begin{center}
\begin{tabular}{ccll}
\toprule
$N_{\mathrm{eff}}$ & Gauge group & $B_{\mathrm{full}}$ & $n^*_{\mathrm{prime}}$ \\
\midrule
4  & $SU(2)\times U(1)$ & $\approx 8.9$   & 17 \\
8  & $SU(3)$ colour     & $\approx 25.2$  & 67 \\
12 & Standard Model     & $\approx 46.3$  & \textbf{137} \\
24 & $SU(5)$ GUT        & $\approx 131.0$ & 467 \\
\bottomrule
\end{tabular}
\end{center}

Different $N_{\mathrm{eff}}$ values produce \emph{different} $n^*$ — the derivation
is non-circular.  Only $N_{\mathrm{eff}} = 12$ (Standard Model mode count) yields
$n^* = 137$.

\subsection{Stress test of correction factor $R \approx 1.114$}

The best algebraic candidate for $R$ is:
\begin{equation}
  R \approx 1 + \alpha(N_{\mathrm{eff}} + \pi + 1/4) \approx 1.1123
  \quad (\text{error } 0.15\%).
  \label{eq:ar:R_candidate}
\end{equation}
\textbf{This formula uses $\alpha$ as input — it is CIRCULAR and is therefore
rejected.}  The value $R \approx 1.114$ is treated as a motivated conjecture [MC]
with no non-circular derivation.  Any paper claiming $\alpha^{-1}_{\mathrm{bare}} = 137$
from this route must acknowledge $R$ as an open problem or use
$B = B_{\mathrm{base}}$ alone (which gives $n^*_{\mathrm{prime}} \neq 137$ for the
naive counting; see validate\_B\_coefficient.py for confirmation).

\subsection{Modular weight coincidence}

The modular weight of $\hat{Z}(\tau) = \vartheta_3^3(\tau)$ is $3/2$, which
equals the exponent in $B_{\mathrm{base}} = N_{\mathrm{eff}}^{3/2}$.
\textbf{Stress test}: is this weight chosen to match the exponent?
No — the weight $3/2$ follows algebraically from $\vartheta_3^3$ having weight
$3 \times 1/2 = 3/2$ where $1/2$ is the weight of $\vartheta_3$ (half-integer
weight from the Jacobi theta function definition).  The exponent $3 = \dim(\mathrm{Im}\,\mathbb{H})$
is fixed by the quaternion algebra.  The coincidence is structural, not fitted.

%──────────────────────────────────────────────────────────────────────────────
\section{Complete Status Table}
\label{sec:ar:status}
%──────────────────────────────────────────────────────────────────────────────

\begin{center}
\renewcommand{\arraystretch}{1.35}
\begin{tabular}{clllll}
\toprule
Step & Result & Status & Circular? & Gap & Source \\
\midrule
1 & $\mathcal{B} = \mathbb{C}\otimes\mathbb{H}$ & CLEAN [L0] & No & — & \S\ref{sec:ar:algebra} \\
2 & $S^1_\psi$, Dirac quantisation & CLEAN [L0] & No & — & \S\ref{sec:ar:compactification} \\
3 & $N_{\mathrm{eff}} = 12$ & CLEAN [L0] & No & — & \S\ref{sec:ar:Neff} \\
4 & $B_0 = 8\pi$ & CLEAN [L1] & No & — & \S\ref{sec:ar:B0} \\
5 & $B_{\mathrm{base}} = N_{\mathrm{eff}}^{3/2}$ & [MC] & No & \textbf{G137-B} & \S\ref{sec:ar:Bbase} \\
6 & $V_{\mathrm{eff}}(n) = n^2 - B\,n\ln n$ & CLEAN [L1] given $B$ & No & — & \S\ref{sec:ar:Veff} \\
7 & $n^* = 137$ (prime attractor) & [L1] conditional & No & via Step~5 & \S\ref{sec:ar:nstar} \\
8 & $\alpha^{-1}_{\mathrm{bare}} = 137$ & CONDITIONAL & No & G137-B & \S\ref{sec:ar:identification} \\
\midrule
\multicolumn{6}{l}{\emph{Not in clean chain (circular inputs):}} \\
9 & $\delta = 0.036$ (two-loop QED) & [CIRC] & Yes & A9 & — \\
10 & $R_\psi$ in SI units & [SE] & Partial & A10 & — \\
11 & $R \approx 1.114$ & [MC]/[CIRC?] & Possibly & — & Stress-tested \\
\bottomrule
\end{tabular}
\end{center}

\begin{resultbox}
\textbf{Summary}: Steps 1--4 and the prime-stability part of Step 7 are proved
with zero free parameters.  Step 5 (Gap G137-B: $k=1$) is the unique remaining gap.
No circular input appears in Steps 1--8.  If Gap G137-B is resolved,
$\alpha^{-1}_{\mathrm{bare}} = 137$ is a zero-parameter result from UBT axioms.
\end{resultbox}

%──────────────────────────────────────────────────────────────────────────────
\section{Gap G137-B: Modular Bootstrap Attack}
\label{sec:ar:bootstrap}
%──────────────────────────────────────────────────────────────────────────────

\begin{gapbox}
\textbf{Gap G137-B}: Prove $k = 1$ (Kac-Moody level) from the UBT CFT on $T^2$.
\end{gapbox}

\subsection{Why the Modular Bootstrap Is the Right Tool}

A complete proof of $k = 1$ requires establishing that the biquaternion field
theory on $T^2$ is a 2D CFT with $SU(2)_1$ (or $U(1)^3$ at level 1) current algebra.
The modular bootstrap approach proceeds:

\begin{enumerate}
  \item Verify that $S[\Theta]$ on $T^2$ is conformally invariant (a 2D CFT).
  \item Compute the 4-point function $\langle\Theta(z_1)\Theta(z_2)\Theta(z_3)\Theta(z_4)\rangle$.
  \item Impose crossing symmetry: the $s$-channel OPE equals the $t$-channel OPE.
  \item Determine which values of $k$ are consistent with the operator spectrum
        dictated by $N_{\mathrm{eff}} = 12$ and partition function
        $\hat{Z}(\tau) = \vartheta_3^3(\tau)$.
\end{enumerate}

\subsection{New Assessment (This Mission)}

The partition function $\hat{Z}(\tau) = \vartheta_3^3(\tau)$ of weight $3/2$ is
consistent with:
\begin{itemize}
  \item Three free compact bosons at the self-dual radius ($k \to \infty$ in the
        single-SU(2) parametrisation, $c = 3$).
  \item Three independent $SU(2)_1$ WZW models ($k = 1$ per Im($\mathbb{H}$) direction,
        $c = 1$ per factor, $c_{\mathrm{total}} = 3$).
\end{itemize}

Both interpretations give the same partition function $\vartheta_3^3(\tau)$ and the
same central charge $c = 3$ (they coincide at the self-dual radius).  The modular
weight alone does not distinguish between them.

\medskip
The modular bootstrap can distinguish them via the operator spectrum:
\begin{itemize}
  \item $SU(2)_1$ (per factor): only primaries $j = 0$ and $j = 1/2$ with
        dimensions $h = 0$ and $h = 1/4$.
  \item Free boson: primary tower at $h = n^2/(4R^2)$ for all integers $n$.
\end{itemize}
Crossing symmetry at the level of the 4-point function will see which operator
content is consistent with the UBT field content.  This is an explicit computation.

\subsection{Time-Box Decision}

\begin{center}
\begin{tabular}{ll}
\toprule
Start date & 2026-04-28 \\
End date & 2026-05-26 (4 weeks) \\
\midrule
Success condition & $k = 1$ forced by crossing symmetry of UBT 4-point function \\
Failure condition & No unique $k$ from bootstrap; route exhausted \\
\midrule
If success & Draft ``$\alpha^{-1}_{\mathrm{bare}} = 137$'' minimal paper \\
If failure & Declare T3\_ALPHA time-boxed; redirect to T1+T2 writing \\
\bottomrule
\end{tabular}
\end{center}

\ifdefined\INCLUDEMODE
\else
  \end{document}
\fi
